% fixture chapter c98 (de); made-up content
\chapter{Kiesel}\label{c98}
\section{Haufen}\label{c98.s1}

\begin{theorem}[Anzahl]\label{c98.thm.count}\claim[proof-omitted]{c98.cl.1}{source-quoted}
  Die Anzahl $n$ erfüllt $n \geq 0$.
  \srcref{fake-notes}{Zeile 5}
  Zwei Haufen zu je einem Kiesel haben zusammen
  \num[c98.cl.4]{2.0}{} Kiesel.
\end{theorem}

\begin{equation}\label{c98.eq.sum}\claim{c98.cl.4}{numerically-checked}
  1 + 1 = 2
\end{equation}

\begin{lemma}[Leer]\label{c98.lem.empty}\claim{c98.cl.3}{derived}
  Der leere Haufen hat Anzahl $0$.
\end{lemma}

\begin{equation}\label{c98.eq.circle}\claim{c98.cl.5}{author-stipulated}
  \bigcirc
\end{equation}

\begin{vertiefung}[id=c98.v1, titel={T}, wozu={W}, voraussetzung={V}, zurueck={Z}]
  \begin{equation}\label{c98.eq.page}\claim{c98.cl.2}{source-backed}
    n \geq 0
  \end{equation}
\end{vertiefung}
